\section{问题分析}

\subsection{问题一分析}
每次测向给出以检测点为顶点、半张角为 \(1\degree\) 的角域约束。多个角域与目标圆域求交后形成凸定位区域，问题转化为约束区域构造、凸集直径计算与最小覆盖判定。

\subsection{问题二分析}
第二检测点既要与首次示向方向形成较大的交会角，又要位于可能接收到信号的区域内，同时控制机器狗移动代价。因此需要建立几何可行域并在其上进行多目标选址。

\subsection{问题三分析}
干扰源频道互异但数量未知，机器狗必须联合决策“去哪里、扫哪些频道、何时复测、何时清除”。可将任务描述为带观测更新的状态机，并用总虚拟时间衡量策略效果。

\subsection{问题四分析}
定向源只在其发射方向两侧各 \(90\degree\) 的范围内可见，单次无信号不能证明频道不存在。策略需增加多方位复测和空间覆盖机制，以降低方向性遮蔽造成的漏检风险。
